\documentclass[11pt]{article}

\usepackage{url}            % simple URL typesetting
\usepackage{authblk}
\usepackage{amsfonts}       % blackboard math symbols
\usepackage{graphicx}
\usepackage{amsmath}
\usepackage{amsthm}
\usepackage[ruled]{algorithm2e}
\usepackage{algorithmic}
\usepackage{subfig}
\usepackage{fullpage}
\usepackage{natbib}

\newtheorem{lemma}{Lemma}

\newtheorem{theorem}{Theorem}
\newtheorem*{theorem*}{Theorem}
\newtheorem{corollary}{Corollary}

\theoremstyle{definition}
\newtheorem{assumption}{Assumption}

\newcommand{\diag}{\text{diag}}

\renewcommand{\a}{{\bf a}}
\renewcommand{\b}{{\bf b}}
\renewcommand{\c}{{\bf c}}
\newcommand{\g}{{\bf g}}
\newcommand{\h}{{\bf h}}
\newcommand{\n}{{\bf n}}
\newcommand{\q}{{\bf q}}
\renewcommand{\r}{{\bf r}}
\newcommand{\s}{{\bf s}}
\newcommand{\w}{{\bf w}}
\renewcommand{\v}{{\bf v}}
\newcommand{\x}{{\bf x}}
\newcommand{\y}{{\bf y}}
\newcommand{\z}{{\bf z}}

\newcommand{\A}{{\bf A}}
\newcommand{\B}{{\bf B}}
\newcommand{\C}{{\bf C}}
\newcommand{\D}{{\bf D}}
\newcommand{\F}{{\bf F}}
\newcommand{\G}{{\bf G}}
\renewcommand{\H}{{\bf H}}
\newcommand{\I}{{\bf I}}
\newcommand{\M}{{\bf M}}
\newcommand{\N}{{\bf N}}
\renewcommand{\P}{{\bf P}}
\newcommand{\Q}{{\bf Q}}
\newcommand{\R}{\mathbb{R}}
\renewcommand{\S}{{\bf S}}
\newcommand{\U}{{\bf U}}
\newcommand{\V}{{\bf V}}
\newcommand{\W}{{\bf W}}
\newcommand{\X}{{\bf X}}
\newcommand{\Y}{{\bf Y}}
\newcommand{\Z}{{\bf Z}}

\newcommand{\EE}{\mathbb{E}}
\newcommand{\RR}{\mathbb{R}}

\newcommand{\calS}{\mathcal{S}}
\newcommand{\calX}{\mathcal{X}}
\newcommand{\calY}{\mathcal{Y}}

\newcommand{\eps}{\varepsilon}

\newcommand{\Gam}{\boldsymbol{\Gamma}}
\newcommand{\Lam}{\boldsymbol{\Lambda}}
\newcommand{\Sig}{\boldsymbol{\Sigma}}

\newcommand{\lb}{\left(}
\newcommand{\rb}{\right)}
\newcommand{\wh}{\widehat}
\newcommand{\wt}{\widetilde}
\newcommand{\ol}{\overline}

\newcommand{\calL}{\mathcal{L}}

\DeclareMathOperator{\tr}{Tr}

\newcommand{\argmax}[1]{\underset{#1}{\operatorname{arg}\,\operatorname{max}}\;}
\newcommand{\argmin}[1]{\underset{#1}{\operatorname{arg}\,\operatorname{min}}\;}
\newcommand{\norm}[1]{\lVert{#1}\rVert}


\usepackage[cmyk]{xcolor}
\newcommand\note[1]{\textcolor{red}{{[#1]}}}

\title{Optimal projection}
\author{David Lipshutz}
\date{\today}

\begin{document}

\maketitle

Let $\x_t$ be a time series and suppose our goal is to project $\x_t$ onto the vector $\v$. Suppose, however, we do not have access to $\x_t$, but only a noisy observation $\y_t=\x_t+\n_t$, where $\n_t$ is isotropic gaussian noise with variance $\alpha^2$. What is the optimal projection of $\y_t$? Consider the optimization problem
\begin{align*}
    \min_{\w} E[|\w^\top\y_t-\v^\top\x_t|^2]&=\min_\w E[|(\w-\v)^\top\x_t+\w^\top\n_t|^2]\\
    &=\min_\w E[|(\w-\v)^\top\x_t|^2]+\alpha^2\w^\top\w\\
    &=\min_\w (\w-\v)^\top\C_{xx}(\w-\v)+\alpha^2\w^\top\w,
\end{align*}
where we have used the fact that $\n_t$ is mean zero isotropic noise independent of $\x_t$. Solving for the optimal $\w$:
\begin{align*}
    \w=(\C_{xx}+\alpha^2\I)^{-1}\C_{xx}\v.
\end{align*}
Note that this is strongly dependent on $\C_{xx}$, which is difficult to interpret in our setting because $\x_t$ is not stationary. If we assume that $\C_{xx}$ is identity, then increasing noise (i.e., $\alpha$) simply scales the optimal project, but does not change the optimal \textit{direction}.

% What if $\n_t$ is adversarial? In this case, the optimization problem becomes:
% \begin{align*}
%     \min_{\w}\max_\n E[|(\w-\v)^\top\x_t+\w^\top\n_t|^2]&=\min_\w\max_\n E[|(\w-\v)^\top\x_t|^2]+2(\w-\v)^\top E[\x_t\n_t^\top]\w+\alpha^2\w^\top E[\n_t\n_t^\top]\w\\
%     &=\min_\w (\w-\v)^\top\C_{xx}(\w-\v)+\alpha^2\w^\top\w,
% \end{align*}

\subsection{RRR}

Suppose we apply Bio-RRR to the time series $\y_t$ (let's focus on the case Bio-RRR corresponds to CCA). Then the optimal projection $\w$ is the left singular vector of the matrix
\begin{align*}
    \C_{yy}^{-1/2}E[\y_t\y_{t+1}^\top]\C_{yy}^{-1/2}&=(\C_{xx}+\alpha\I)^{-1/2}E[\x_t\x_{t+1}^\top+\n_t\n_{t+1}^\top](\C_{xx}+\alpha\I)^{-1/2}\\
    &=(\C_{xx}+\alpha\I)^{-1/2}(\A^\top+\alpha\S)(\C_{xx}+\alpha\I)^{-1/2}.
\end{align*}
Here $\S$ is the matrix of off-diagonal ones defined in section III.B of the main text. How does the left-singular vector $\w$ change with $\alpha$? When $\alpha=0$, then optimal projection is the left singular vector of $\C_{xx}^{-1/2}\A^\top\C_{xx}^{-1/2}$ (right singular vector of $\C_{xx}^{-1/2}\A\C_{xx}^{-1/2}$). As $\alpha\to\infty$, the optimal projection converges to the left singular vector of $\S$.

\end{document}